
% -------------------------------------------------------
%  Common Styles and Formattings
% -------------------------------------------------------


\usepackage{amssymb,amsmath}
\usepackage[colorlinks,linkcolor=blue,citecolor=blue]{hyperref}
\usepackage{geometry}
\usepackage{pdfpages}
\usepackage{physics}
\usepackage{pythonhighlight}
\usepackage{fancyhdr}
\usepackage{afterpage}
\usepackage{efbox}



% 
\usepackage{listings}
\usepackage{mathtools,amsthm}
\usepackage{setspace}
\usepackage[chapter]{minted} 
\usepackage[ruled,linesnumbered,algochapter]{algorithm2e}
\usepackage{enumitem}
\usepackage{booktabs}
\usepackage[table]{xcolor}
\usepackage{tabularx}
\usepackage{multirow}
\usepackage{comment}
\usepackage{subfig}
\usepackage{graphicx}
\usepackage{tikz}
\usetikzlibrary{arrows.meta, positioning, shapes.geometric, fit, calc, patterns, backgrounds, angles, quotes, decorations.markings}
\usepackage{forest}
\usepackage{float}
\usepackage{longtable}

% سبک مشترک نمودارهای درختی زیرگروه شبیه‌سازی (رشد افقی از چپ به راست)
\forestset{
	qsimtree/.style={
		for tree={
			draw,
			rounded corners,
			grow'=east,
			align=left,
			font=\footnotesize,
			edge={-latex},
			parent anchor=east,
			child anchor=west,
			anchor=west,
			l sep=8mm,
			s sep=1.6mm,
			minimum height=6mm,
			inner xsep=5pt,
			inner ysep=3pt,
			edge path={
				\noexpand\path[\forestoption{edge}]
				(!u.parent anchor) -- +(4mm,0) |- (.child anchor)\forestoption{edge label};
			}
		}
	}
}
%

% 
\usepackage{array}
%


\usepackage[extrafootnotefeatures]{xepersian}

\usepackage[bottom]{footmisc}

% -------------------- Page Layout --------------------

\efboxsetup{linewidth=2pt}

%\newgeometry{top=3cm,right=3cm,left=2.5cm,bottom=3cm,footskip=1.25cm}
\newgeometry{margin=1in,bottom=1.1in,footskip=.4in}

\renewcommand{\baselinestretch}{1.4}
\linespread{1.6}
\setlength{\parskip}{0.45em}


\fancyhf{}
\renewcommand{\headrulewidth}{0pt}
\cfoot{\thepage}
\pagestyle{fancy}   


% -------------------- Default Font --------------------

\settextfont[
Scale=1.09,
Extension=.ttf, 
Path=styles/fonts/,
BoldFont=XB NiloofarBd,
ItalicFont=XB NiloofarIt,
BoldItalicFont=XB NiloofarBdIt
]{XB Niloofar}

%\setdigitfont[
%Scale=1.09,
%Extension=.ttf, 
%Path=styles/fonts/,
%BoldFont=XB NiloofarBd,
%ItalicFont=XB NiloofarIt,
%BoldItalicFont=XB NiloofarBdIt
%]{XB Niloofar}


% -------------------- Sharif Font --------------------

\newcommand{\shariffont}{
\settextfont[
Scale=1.4,
Extension=.ttf, 
Path=styles/fonts/,
BoldFont=Sharif1.3-SemiBold,
]{Sharif1.3-Regular}

\setdigitfont[
Scale=1.4,
Extension=.ttf, 
Path=styles/fonts/,
BoldFont=Sharif1.3-SemiBold,
]{Sharif1.3-Regular}
}


% -------------------- Titles --------------------


%\renewcommand{\listfigurename}{فهرست شکل‌ها}
%\renewcommand{\listtablename}{فهرست جدول‌ها}
%\renewcommand{\bibname}{\rl{{مراجع}\hfill}} 


% -------------------- Commands --------------------

\newcommand{\ft}[1]{\LTRfootnote{ #1}}

% ------------------------------ Images and Figures --------------------------

\graphicspath{{figs/}{figures/}{front/template/}}

%\newcommand{\centerimg}[1]{
%%	\vspace{1em}
%	\begin{figure}[!htbp]
%		\begin{center}
%			\efbox{\includegraphics[width=\textwidth]{figs/#1}}
%		\end{center}
%	\end{figure}
%}

% ------------------------------ NK --------------------------

\renewcommand\arraystretch{0.6}  % for matrix size

\renewcommand*{\thepart}{\arabic{part}}

\makeatletter
\def\@part[#1]#2{%
	\ifnum \c@secnumdepth > -2\relax
	\refstepcounter{part}%
	% *** THIS is the important line for the TOC: ***
	\addcontentsline{toc}{part}{بخش \thepart \quad #1}%
	\else
	\addcontentsline{toc}{part}{#1}%
	\fi
	\markboth{}{}%
	\addvspace{50\p@}%
	\begingroup
	\centering
	\interlinepenalty\@M
	\normalfont
	\Huge \bfseries \partname\nobreakspace\thepart
	\par
	\vskip 20\p@
	{\Huge \bfseries #2\par}%
	\endgroup
	\@endpart}
\makeatother

\makeatletter
	\@addtoreset{chapter}{part}
\makeatother

% --------------------------------------------------------

\definecolor{codebackcolor}{rgb}{0.95,0.95,0.95}

\renewcommand{\theFancyVerbLine}{
	\sffamily
	\textcolor[rgb]{0.5,0.5,1.0}{
		\tiny\arabic{FancyVerbLine}
	}
}

\newcommand{\codespacing}{1.6}

% #1 language #2 file #3 label
\newcommand{\code}[3]{
\begin{latin}
	\begin{spacing}{\codespacing}
		\inputminted[breaklines=true,fontsize=\footnotesize,linenos=true,frame=leftline,bgcolor=codebackcolor]{#1}{#2}
	\end{spacing}
	\vspace{2mm}
	\label{#3}
\end{latin}
}



\makeatletter
\def\thm@space@setup{\thm@preskip=0pt
	\thm@postskip=0pt}
\makeatother
% Theorem, Lemma, Corollary, Proposition, Conjecture, Criterion, Assertion
\theoremstyle{plain} % italic text, extra space above and below.
\newtheorem{theorem}{\bf قضیه}[chapter]
\newtheorem{proposition}{\bf گزاره}[chapter]
\newtheorem{lemma}{\bf لم}[chapter]
\newtheorem{corollary}{\bf نتیجه فرعی}[chapter]
\newtheorem{assumption}{\bf فرض}[chapter]
\newtheorem{conjecture}{\bf حدس}[chapter]
% Definition, Condition, Problem, Example, Exercise, Algorithm, Question, Axiom, Property, Assumption, Hypothesis
\theoremstyle{definition} % upright text, extra space above and below
\newtheorem{definition}{\bf تعریف}[chapter]
\newtheorem{question}{\bf سوال}[chapter]
\newtheorem{oepn_question}{\bf سوال باز}[chapter]
\newtheorem{example}{\bf مثال}[chapter]
\newtheorem{problem}{\bf مسئله}[chapter]
\newtheorem{algorithm_env}{\bf الگوریتم}[chapter] % algorithmm is a perdifined environment, so we named it algoritm_env
%\newtheorem{assumption}{\bf فرض}[chapter]
% Remark, Note, Notation, Claim, Summary, Acknowledgement, Case, Conclusion
\theoremstyle{remark} % Upright text, no extra space above or below.
\newtheorem{claim}{\bf ادعا}[chapter]
\newtheorem{note}{\bf توجه}[chapter]
\newtheorem{remark}{\bf نکته}[chapter]

\DeclareMathOperator{\module}{mod} % Default \mod has extra spaces.

\newenvironment{nospacecenter}
{\parskip=0pt\par\nopagebreak\centering}
{\par\noindent\ignorespacesafterend}

\newenvironment{mylisting}{\captionsetup{type=listing}}{}

%----------------------------------------------------

\newcommand{\centerimg}[2]
{
	\vspace{1em}
	\begin{center}\includegraphics[width=#2]{figs/#1}\end{center}
	\vspace{-1em}
}